\documentclass[UTF8,a4paper,12pt]{ctexart}

\usepackage{geometry}
\usepackage{booktabs}
\usepackage{longtable}
\usepackage{array}
\usepackage{graphicx}
\usepackage{enumitem}
\usepackage{hyperref}
\usepackage{xcolor}
\usepackage{amsmath}
\usepackage{float}

\geometry{left=2.5cm,right=2.5cm,top=2.5cm,bottom=2.5cm}

\hypersetup{
	colorlinks=true,
	linkcolor=blue,
	urlcolor=blue,
	citecolor=blue
}

\title{\textbf{第 4 周实验计划：实现两阶段检索 Coarse-to-Fine}}
\author{}
\date{2026.4.6 -- 2026.4.12}

\begin{document}
	
	\maketitle
	
	\section{本周定位与总体目标}
	
	本周进入真正的 budgeted retrieval 设计阶段。与前几周不同，本周的目标不是继续增加一个 baseline，而是建立一个可控制候选规模的两阶段检索框架，并首次系统分析检索质量与计算成本之间的关系。
	
	本周的核心思想是：先使用低成本的 coarse retriever 从全库中召回有限数量的候选页面，再仅在这些候选页面上执行 full late interaction reranking。这样可以避免对整个语料库进行昂贵的 full multi-vector matching，从而将 multi-vector retrieval 转化为一个可预算、可调节、可分析的检索系统。
	
	本周重点关注如下问题：
	
	\begin{itemize}[leftmargin=2em]
		\item 在不同候选预算下，coarse-to-fine 检索能保留多少 full multi-vector 的效果？
		\item reranking 成本是否随候选规模近似线性增长？
		\item 哪一个候选规模能够取得最好的质量--成本折中？
		\item 质量损失主要来自 coarse 阶段漏召回，还是来自 fine 阶段排序能力不足？
	\end{itemize}
	
	\section{本周预期成果}
	
	本周结束时，应至少获得以下成果：
	
	\begin{table}[H]
		\centering
		\caption{第 4 周预期成果}
		\begin{tabular}{p{7cm}p{6cm}}
			\toprule
			\textbf{成果} & \textbf{目标状态} \\
			\midrule
			coarse retrieval + fine rerank 管线 & 已完整跑通 \\
			不同 N 设置下的结果 & 已形成系统对比 \\
			coarse-to-fine 与 full MV 的对照结果 & 已可比较 \\
			平均时延与重排成本记录 & 已完整统计 \\
			第一张质量--成本曲线 & 已绘制初版 \\
			最具性价比的 N 候选区间 & 已有初步结论 \\
			\bottomrule
		\end{tabular}
	\end{table}
	
	本周关键词为：\textbf{候选预算、精排效率、质量--成本曲线、Pareto 线索}。
	
	候选规模设置建议为：
	
	$$
	N \in \{10, 20, 50, 100\}
	$$
	
	若语料库规模较小，也可以调整为：
	
	$$
	N \in \{5, 10, 20, 50\}
	$$
	
	\section{方法设计：Coarse-to-Fine 两阶段检索}
	
	\subsection{标准流程}
	
	本周固定采用如下两阶段检索流程：
	
	\begin{enumerate}[leftmargin=2em]
		\item 使用低成本 coarse retriever 从全库召回 top-N 页面；
		\item 对 top-N 候选页面执行 full late interaction reranking；
		\item 输出最终 top-k 排序结果；
		\item 评测检索效果、时延和重排成本；
		\item 与 full multi-vector retrieval 进行对照分析。
	\end{enumerate}
	
	形式化地，两阶段检索可以表示为：
	
	$$
	\mathcal{C}_N(q) = \operatorname{TopN}_{p \in \mathcal{D}} s_c(q,p)
	$$
	
	其中，$$q$$ 表示 query，$$p$$ 表示页面，$$\mathcal{D}$$ 表示全库页面集合，$$s_c(q,p)$$ 表示 coarse retriever 的相似度分数，$$\mathcal{C}_N(q)$$ 表示 query 对应的 top-N 候选页面集合。
	
	随后，在候选集合内执行 late interaction reranking：
	
	$$
	s_f(q,p) = \operatorname{LateInteraction}(Q, P), \quad p \in \mathcal{C}_N(q)
	$$
	
	最终输出：
	
	$$
	\operatorname{Rank}(q) = \operatorname{Sort}_{p \in \mathcal{C}_N(q)} s_f(q,p)
	$$
	
	该设计的关键优势在于：full late interaction 不再作用于整个语料库，而是仅作用于经过 coarse 阶段筛选后的有限候选集合，从而显著降低 reranking 成本。
	
	\subsection{Coarse Retriever 选择}
	
	本周建议优先使用 Week 2 已经跑通的 single-vector visual retrieval 作为 coarse retriever。
	
	选择理由如下：
	
	\begin{itemize}[leftmargin=2em]
		\item 已有稳定实现，工程风险较低；
		\item 已有 baseline 结果，便于横向比较；
		\item 可自然构成 single-vector coarse retrieval 与 multi-vector fine reranking 的两阶段框架；
		\item 有助于清晰分析 coarse 阶段召回能力对最终结果的影响。
	\end{itemize}
	
	\subsection{成本记录指标}
	
	对每个候选规模 N，需要同步记录如下成本指标：
	
	\begin{table}[H]
		\centering
		\caption{成本记录指标}
		\begin{tabular}{p{5cm}p{8cm}}
			\toprule
			\textbf{指标} & \textbf{说明} \\
			\midrule
			coarse retrieval time & 粗召回阶段耗时 \\
			rerank time & late interaction 精排耗时 \\
			total latency & 单次检索总耗时 \\
			per-query latency & 平均每条 query 的耗时 \\
			top-N candidate size & coarse 阶段输出候选数量 \\
			final rerank size & 实际进入 reranking 的候选数量 \\
			memory / GPU memory & 内存或显存占用，若方便则记录 \\
			\bottomrule
		\end{tabular}
	\end{table}
	
	\section{六天任务安排}
	
	\subsection{第 1 天：明确两阶段检索方案并确定 coarse retriever}
	
	\subsubsection*{目标}
	
	搭建 coarse-to-fine 检索框架的整体骨架，明确 coarse 阶段、fine 阶段、候选规模和成本记录方式。
	
	\subsubsection*{主要任务}
	
	\begin{enumerate}[leftmargin=2em]
		\item 明确两阶段流程：
		\begin{itemize}
			\item coarse retriever 从全库召回 top-N 页面；
			\item full late interaction 仅对 top-N 候选页面精排；
			\item 输出最终 top-k；
			\item 记录效果和时延。
		\end{itemize}
		
		\item 选定 coarse retriever：
		\begin{itemize}
			\item 推荐使用 Week 2 的 single-vector visual retrieval；
			\item 保证与 Week 3 的 full MV reranker 可直接衔接。
		\end{itemize}
		
		\item 固定候选规模：
		
		$$
		N \in \{10, 20, 50, 100\}
		$$
		
		\item 建立实验日志模板，确保每个 N 都记录效果与成本。
	\end{enumerate}
	
	\subsubsection*{当日产出}
	
	\begin{itemize}[leftmargin=2em]
		\item \texttt{coarse\_to\_fine\_design.md}
		\item N 配置表
		\item coarse retriever 选择说明
		\item 本周实验日志模板
	\end{itemize}
	
	\subsubsection*{验收标准}
	
	第 1 天结束时，应明确以下内容：
	
	\begin{itemize}[leftmargin=2em]
		\item coarse 阶段使用什么方法；
		\item fine 阶段如何接入 full MV；
		\item 本周测试哪些 N；
		\item 每个 N 需要记录哪些指标。
	\end{itemize}
	
	\subsection{第 2 天：实现 coarse retrieval 与 candidate pipeline}
	
	\subsubsection*{目标}
	
	完成 coarse retrieval，并为每个 query 输出不同候选深度下的 top-N 候选页面集合。
	
	\subsubsection*{主要任务}
	
	\begin{enumerate}[leftmargin=2em]
		\item 调用 single-vector coarse retriever，从全库召回候选页面。
		\item 一次性输出多个候选深度：
		\begin{itemize}
			\item top-10；
			\item top-20；
			\item top-50；
			\item top-100。
		\end{itemize}
		\item 保存统一候选文件 \texttt{coarse\_candidates.tsv}，字段包括 query 编号、候选页面编号、coarse 排名和 coarse 得分。
	\end{enumerate}
	
	\begin{table}[H]
		\centering
		\caption{Coarse Candidate 文件字段}
		\begin{tabular}{p{5cm}p{8cm}}
			\toprule
			\textbf{字段} & \textbf{说明} \\
			\midrule
			query\_id & 查询编号 \\
			candidate\_page\_id & 候选页面编号 \\
			coarse\_rank & coarse 阶段排名 \\
			coarse\_score & coarse 阶段得分 \\
			\bottomrule
		\end{tabular}
	\end{table}
	
	\subsubsection*{Coarse Recall 统计}
	
	本日重点不是最终排序，而是检查 coarse 阶段是否足够保 recall。建议统计：
	
	\begin{itemize}[leftmargin=2em]
		\item coarse Recall@10；
		\item coarse Recall@20；
		\item coarse Recall@50；
		\item coarse Recall@100。
	\end{itemize}
	
	若 gold page 在 coarse 阶段没有进入 top-N，那么后续 fine reranking 无法恢复该结果。因此，coarse recall 是两阶段检索系统的上限约束。
	
	\subsubsection*{当日产出}
	
	\begin{itemize}[leftmargin=2em]
		\item \texttt{coarse\_candidates.tsv}
		\item coarse 阶段召回统计表
		\item 粗召回失败样例记录
	\end{itemize}
	
	\subsubsection*{验收标准}
	
	第 2 天结束时，应确认：
	
	\begin{itemize}[leftmargin=2em]
		\item coarse retriever 是否能稳定召回相关页面；
		\item 哪个 N 起能够避免明显 recall 损失；
		\item 哪些 query 在 coarse 阶段已经丢失答案。
	\end{itemize}
	
	\subsection{第 3 天：在 top-N 候选上接入 full late interaction 精排}
	
	\subsubsection*{目标}
	
	完成 coarse-to-fine 管线的 second-stage reranking，让 full MV 只在候选集合内执行 late interaction。
	
	\subsubsection*{主要任务}
	
	\begin{enumerate}[leftmargin=2em]
		\item 设计 rerank 输入接口：
		\begin{itemize}
			\item 输入文件为 \texttt{coarse\_candidates.tsv}；
			\item 每个 query 对应 top-N 页面；
			\item query 表示与 Week 3 保持一致；
			\item page multi-vector embeddings 复用 Week 3 结果。
		\end{itemize}
		
		\item 对每个 N 独立输出 run 文件：
		\begin{itemize}
			\item \texttt{c2f\_N10\_run.tsv}
			\item \texttt{c2f\_N20\_run.tsv}
			\item \texttt{c2f\_N50\_run.tsv}
			\item \texttt{c2f\_N100\_run.tsv}
		\end{itemize}
		
		\item 先在小子集上验证：
		\begin{itemize}
			\item query 是否拿到正确候选；
			\item rerank 是否只发生在候选集合内；
			\item 最终 top-k 输出是否合理。
		\end{itemize}
		
		\item 检查排序逻辑：
		\begin{itemize}
			\item coarse 排名较低但相关的页面，是否能被 fine reranker 提升；
			\item 视觉相似但语义不相关的页面，是否被压低。
		\end{itemize}
	\end{enumerate}
	
	\subsubsection*{当日产出}
	
	\begin{itemize}[leftmargin=2em]
		\item 各个 N 对应的 rerank run 文件
		\item coarse-to-fine 管线脚本
		\item 子集验证日志
	\end{itemize}
	
	\subsubsection*{验收标准}
	
	第 3 天结束时，应拥有完整的两阶段系统：
	
	$$
	\text{coarse retrieval} \rightarrow \text{candidate set} \rightarrow \text{late interaction rerank} \rightarrow \text{top-k output}
	$$
	
	\subsection{第 4 天：系统跑出不同 N 下的效果与时延结果}
	
	\subsubsection*{目标}
	
	在正式实验集上运行不同 N 设置，获得效果与时延对比结果。
	
	\subsubsection*{主要任务}
	
	\begin{enumerate}[leftmargin=2em]
		\item 对每个 N 运行完整评测；
		\item 记录 Recall@1、Recall@5、Recall@10、MRR@10、nDCG@10；
		\item 记录 coarse time、rerank time、total latency、per-query latency；
		\item 与 full MV 结果进行对照。
	\end{enumerate}
	
	\subsubsection*{核心结果表}
	
	\begin{table}[H]
		\centering
		\caption{Full MV 与 Coarse-to-Fine 检索结果对比}
		\begin{tabular}{lcccc}
			\toprule
			\textbf{Method} & \textbf{Recall@10} & \textbf{MRR@10} & \textbf{nDCG@10} & \textbf{Latency} \\
			\midrule
			Full Multi-vector & -- & -- & -- & -- \\
			C2F N=10 & -- & -- & -- & -- \\
			C2F N=20 & -- & -- & -- & -- \\
			C2F N=50 & -- & -- & -- & -- \\
			C2F N=100 & -- & -- & -- & -- \\
			\bottomrule
		\end{tabular}
	\end{table}
	
	\subsubsection*{分析重点}
	
	第 4 天重点关注趋势，而不是单个绝对数值：
	
	\begin{itemize}[leftmargin=2em]
		\item 小 N 是否已经接近 full MV；
		\item N 增大后效果是否趋于饱和；
		\item rerank time 是否随 N 线性或近线性增长；
		\item coarse-to-fine 的总时延是否显著低于 full MV。
	\end{itemize}
	
	\subsubsection*{当日产出}
	
	\begin{itemize}[leftmargin=2em]
		\item 所有 N 对应的结果文件
		\item \texttt{c2f\_summary.xlsx}
		\item 初版效果--时延对比表
	\end{itemize}
	
	\subsubsection*{验收标准}
	
	第 4 天结束时，应首次清楚看到：小候选集是否已经能够逼近 full multi-vector retrieval。
	
	\subsection{第 5 天：绘制第一张质量--成本曲线并找出最佳 N}
	
	\subsubsection*{目标}
	
	将第 4 天的实验结果转化为论文图，绘制 coarse-to-fine 的质量--成本曲线，并判断最具性价比的候选规模。
	
	\subsubsection*{主要任务}
	
	\begin{enumerate}[leftmargin=2em]
		\item 绘制质量--成本图：
		\begin{itemize}
			\item 横轴：Latency、rerank cost 或 total cost；
			\item 纵轴：Recall@10 或 nDCG@10。
		\end{itemize}
		
		\item 标记关键点：
		\begin{itemize}
			\item full MV 点；
			\item single-vector 点；
			\item best trade-off 点；
			\item 最小延迟点；
			\item 接近饱和点。
		\end{itemize}
		
		\item 分析最佳 N：
		\begin{itemize}
			\item 效果损失是否较小；
			\item 时延节省是否明显；
			\item rerank 成本是否显著下降。
		\end{itemize}
		
		\item 按 query type 进行补充分析：
		\begin{itemize}
			\item 哪类 query 对小 N 最敏感；
			\item 哪类 query 在 coarse 阶段最容易丢失证据页；
			\item 哪类 query 即使 N 很小也能稳定命中。
		\end{itemize}
	\end{enumerate}
	
	\subsubsection*{质量--成本图占位}
	
	\begin{figure}[H]
		\centering
		\fbox{
			\begin{minipage}[c][6cm][c]{0.75\linewidth}
				\centering
				质量--成本曲线图占位\\[0.5em]
				横轴：Latency / Rerank Cost / Total Cost\\
				纵轴：Recall@10 / nDCG@10\\[0.5em]
				后续可替换为 \texttt{quality\_cost\_curve.pdf}
			\end{minipage}
		}
		\caption{Coarse-to-Fine 检索的质量--成本曲线。横轴表示检索成本或时延，纵轴表示 Recall@10 或 nDCG@10。图中应标注 full MV、single-vector 以及不同 N 下的 C2F 结果点。}
		\label{fig:quality_cost_curve}
	\end{figure}
	
	\subsubsection*{预期结论模板}
	
	可形成如下结论：
	
	\begin{quote}
		当 N=20 或 N=50 时，coarse-to-fine 已能保留 full multi-vector retrieval 的大部分检索效果，同时显著降低 reranking 成本，说明适当的候选预算能够带来较优的质量--效率折中。
	\end{quote}
	
	\subsubsection*{当日产出}
	
	\begin{itemize}[leftmargin=2em]
		\item 第一张质量--成本曲线图
		\item 最佳 N 分析笔记
		\item 图注草稿
	\end{itemize}
	
	\subsubsection*{验收标准}
	
	第 5 天结束时，应能够明确回答：
	
	\begin{itemize}[leftmargin=2em]
		\item 哪个 N 最有性价比；
		\item 质量损失主要发生在哪个阶段；
		\item coarse-to-fine 是否值得继续推进。
	\end{itemize}
	
	\subsection{第 6 天：整理结果、分析失败案例并写本周实验小结}
	
	\subsubsection*{目标}
	
	将本周工作从“方法能跑”整理为“论文能写”，完成结果归纳、失败模式分析与 Week 5 的设计过渡。
	
	\subsubsection*{典型案例整理}
	
	建议整理三类典型案例。
	
	\paragraph{A. C2F 明显逼近 full MV 的案例}
	
	这类案例说明：
	
	\begin{itemize}[leftmargin=2em]
		\item coarse 阶段已经成功召回 relevant page；
		\item fine 阶段将 relevant page 排到前列；
		\item 候选集较小但最终效果很好。
	\end{itemize}
	
	\paragraph{B. 小 N 失败、大 N 成功的案例}
	
	这类案例说明：
	
	\begin{itemize}[leftmargin=2em]
		\item 问题主要不在 rerank；
		\item 主要失败原因是 coarse 阶段在小 N 下漏召回；
		\item 增大候选预算后，fine reranker 能恢复正确排序。
	\end{itemize}
	
	\paragraph{C. 即使大 N 仍失败的案例}
	
	这类案例说明可能存在更深层问题：
	
	\begin{itemize}[leftmargin=2em]
		\item query 表示不足；
		\item 页面局部视觉特征不足；
		\item 单页检索无法解决跨页证据需求；
		\item 数据标注或评测目标存在歧义。
	\end{itemize}
	
	\subsubsection*{周报内容}
	
	本周周报建议包括：
	
	\begin{itemize}[leftmargin=2em]
		\item coarse-to-fine 方法实现情况；
		\item 不同 N 的总体结果；
		\item 时延与重排成本比较；
		\item 最佳 N 的初步判断；
		\item 本周发现的失败模式；
		\item 对 Week 5 patch/token selection 的启示。
	\end{itemize}
	
	\subsubsection*{Week 5 过渡逻辑}
	
	本周 coarse-to-fine 主要解决的是“候选页面数太多”的问题。自然地，Week 5 可以进一步解决“每页 token 太多”的问题。
	
	因此，论文逻辑可以组织为：
	
	\begin{quote}
		Coarse-to-fine retrieval reduces the number of pages that require expensive multi-vector matching, while patch or token selection further reduces the number of visual tokens involved in each page-level late interaction. Together, these two strategies form a budgeted multi-vector retrieval framework.
	\end{quote}
	
	\subsubsection*{当日产出}
	
	\begin{itemize}[leftmargin=2em]
		\item \texttt{weekly\_report\_week4.md}
		\item 质量--成本曲线图初稿
		\item 典型案例截图
		\item Week 5 设计笔记
	\end{itemize}
	
	\subsubsection*{验收标准}
	
	第 6 天结束时，应能够写出如下论文级别结论：
	
	\begin{quote}
		通过将 full late interaction 限制在 top-N 候选页面上，两阶段检索能够在显著降低精排成本的同时保留大部分检索质量，证明 coarse-to-fine 是预算化 multi-vector retrieval 的有效组成部分。
	\end{quote}
	
	\section{实验记录模板}
	
	\subsection{不同 N 的效果记录}
	
	\begin{table}[H]
		\centering
		\caption{不同候选规模下的检索效果}
		\begin{tabular}{lccccc}
			\toprule
			\textbf{Method} & \textbf{Recall@1} & \textbf{Recall@5} & \textbf{Recall@10} & \textbf{MRR@10} & \textbf{nDCG@10} \\
			\midrule
			Single-vector & -- & -- & -- & -- & -- \\
			Full MV & -- & -- & -- & -- & -- \\
			C2F N=10 & -- & -- & -- & -- & -- \\
			C2F N=20 & -- & -- & -- & -- & -- \\
			C2F N=50 & -- & -- & -- & -- & -- \\
			C2F N=100 & -- & -- & -- & -- & -- \\
			\bottomrule
		\end{tabular}
	\end{table}
	
	\subsection{不同 N 的成本记录}
	
	\begin{table}[H]
		\centering
		\caption{不同候选规模下的检索成本}
		\begin{tabular}{lcccc}
			\toprule
			\textbf{Method} & \textbf{Coarse Time} & \textbf{Rerank Time} & \textbf{Total Latency} & \textbf{Per-query Latency} \\
			\midrule
			Single-vector & -- & -- & -- & -- \\
			Full MV & -- & -- & -- & -- \\
			C2F N=10 & -- & -- & -- & -- \\
			C2F N=20 & -- & -- & -- & -- \\
			C2F N=50 & -- & -- & -- & -- \\
			C2F N=100 & -- & -- & -- & -- \\
			\bottomrule
		\end{tabular}
	\end{table}
	
	\subsection{Coarse Recall 记录}
	
	\begin{table}[H]
		\centering
		\caption{Coarse 阶段召回能力}
		\begin{tabular}{lcccc}
			\toprule
			\textbf{Retriever} & \textbf{Recall@10} & \textbf{Recall@20} & \textbf{Recall@50} & \textbf{Recall@100} \\
			\midrule
			Single-vector coarse retriever & -- & -- & -- & -- \\
			\bottomrule
		\end{tabular}
	\end{table}
	
	\section{本周核心分析问题}
	
	本周结果分析应围绕以下几个问题展开：
	
	\begin{enumerate}[leftmargin=2em]
		\item Coarse-to-fine 在多大程度上接近 full MV？
		\item 当 N 增大时，效果提升是否迅速饱和？
		\item rerank time 是否与 N 近似线性相关？
		\item 最优 N 是由效果决定，还是由时延预算决定？
		\item 失败案例主要是 coarse 阶段漏召回，还是 fine 阶段排序失败？
		\item 本周结论如何支持 Week 5 的 patch/token selection 设计？
	\end{enumerate}
	
	\section{本周总结结论}
	
	本周的目标是实现并验证一个可控制候选规模的 coarse-to-fine retrieval 框架。该框架首先使用低成本 single-vector retriever 从全库中召回 top-N 候选页面，然后仅在候选集合内执行 full late interaction reranking。通过系统比较不同 N 下的检索效果、时延和 reranking 成本，可以得到第一条质量--成本曲线，并初步判断最具性价比的候选预算。
	
	预期论文级结论如下：
	
	\begin{quote}
		通过将 full late interaction 限制在 top-N 候选页面上，两阶段检索能够在显著降低精排成本的同时保留大部分检索质量。实验结果表明，适当的候选预算可以形成有效的质量--效率折中，证明 coarse-to-fine retrieval 是预算化 multi-vector retrieval 框架中的关键组成部分。
	\end{quote}
	
	进一步地，本周结果将自然引出下一阶段研究：如果 coarse-to-fine 主要减少需要精排的页面数量，那么 Week 5 的 patch/token selection 将进一步减少每个页面内部参与 late interaction 的 token 数量，从而形成更加完整的 budgeted multi-vector retrieval 方法。
	
\end{document}
